% ID: SEL-D0028-OA-E01
% ESTADO: validada. S029, cotejada con examen y criterios.
\textbf{a)} Las cotas son $0\le y\le\min(2x-2,(x+2)/2,15-3x)$.
Para $y=0$ se obtiene $1\le x\le5$. Las dos primeras fronteras
se cortan en $(2,2)$, y la segunda y tercera en $(4,3)$.
La otra intersección de fronteras superiores supera la cota $(x+2)/2$.
El recinto es el cuadrilátero siguiente:
\begin{center}
\setlength{\unitlength}{8mm}
\begin{picture}(8,5)(-1,-1)
\qbezier(0,0)(3,0)(6,0)\put(6.2,0){$x$}
\qbezier(0,0)(0,1.8)(0,3.6)\put(0,3.7){$y$}
\qbezier(1,0)(1.5,1)(2,2)
\qbezier(2,2)(3,2.5)(4,3)
\qbezier(4,3)(4.5,1.5)(5,0)
\put(1,-0.4){\makebox(0,0)[t]{$(1,0)$}}
\put(2,2.2){\makebox(0,0)[b]{$(2,2)$}}
\put(4,3.2){\makebox(0,0)[b]{$(4,3)$}}
\put(5,-0.4){\makebox(0,0)[t]{$(5,0)$}}
\end{picture}
\end{center}
\textbf{b)} En $(3,3)$ se tiene $-x+2y=3>2$, así que queda fuera.

\textbf{c)} Una función lineal alcanza sus extremos en vértices.
En $(1,0)$, $(2,2)$, $(4,3)$ y $(5,0)$ los valores son $3,2,6,15$.
El mínimo es $\boxed{2}$ sólo en $(2,2)$ y el máximo
$\boxed{15}$ sólo en $(5,0)$.
